\chapter{Testing and Evaluation}

Where the previous chapter quantified the \emph{model's} predictive performance,
this chapter evaluates the \emph{application} that delivers it. It reports the
automated unit-test suite, functional test cases over the end-to-end pipeline,
boundary and edge-case behaviour, on-device performance, usability testing with
representative users, and a set of use-case scenarios that exercise the full
system. Together these confirm that the system is not only accurate (Results
chapter) but also robust, responsive, and usable on a real device.

% ============================================================
\section{Testing Methodology}
% ============================================================

\subsection{Testing Objectives}

The objectives of application testing were to verify that (i)~the core
model, security, and analytics logic behaves correctly in isolation; (ii)~the
three-stage inference pipeline produces the correct outcome end-to-end for each
class of input, including the normal early-exit and the borderline abstention
paths; (iii)~the application degrades gracefully on extreme or malformed input
rather than crashing; (iv)~inference completes within an acceptable latency and
storage budget on a real device; and (v)~the interface is usable without
instruction by users of differing backgrounds.

\subsection{Testing Strategy}

Testing combined automated and manual methods. Automated \emph{unit tests}
(\texttt{flutter test}) cover the deterministic logic layers. Manual
\emph{functional} and \emph{boundary} testing exercised the full pipeline on a
physical Android device across normal, diseased, borderline, and adversarial
inputs. \emph{Performance} testing measured per-stage inference latency and
model footprint on-device, and \emph{usability} testing observed three
representative participants completing unguided tasks.

% ============================================================
\section{Unit Testing}
% ============================================================

The Dart unit suite (\texttt{flutter test}) contains \textbf{23 tests across five
files}; all pass. Coverage spans the model layer, the security layer, the
progression analytics, and the body-part selection UI
(Table~\ref{tab:unit_tests}).

\begin{table}[H]
\centering
\caption{Unit test coverage (\texttt{flutter test}: 23 passed, 0 failed).}
\label{tab:unit_tests}
\begin{tabular}{|p{3.7cm}|c|p{7.3cm}|}
\hline
\textbf{Test file} & \textbf{\#} & \textbf{Coverage} \\ \hline
\texttt{appointment\_test.dart} & 4 & Appointment model \texttt{toMap}/\texttt{fromMap} round-trip; default status; \texttt{copyWith} isolation; missing-field tolerance \\ \hline
\texttt{password\_hasher\_test.dart} & 6 & PBKDF2-SHA256 hash format; correct/incorrect verification; random-salt uniqueness; legacy-plaintext detection; malformed-hash rejection \\ \hline
\texttt{progression\_utils\_test.dart} & 5 & Group-by-body-part with chronological sort; \texttt{Unspecified} bucket; \texttt{trackableBodyParts} ordering/filtering; \texttt{confidenceTrend} normal-vs-disease mapping \\ \hline
\texttt{widget\_test.dart} & 2 & User model round-trip; \texttt{fromMap} defaults \\ \hline
\texttt{body\_part\_render\_test.dart} & 6 & Body-part screen: header + empty-state hint; hotspot selection; re-tap clears; single-selection replacement; front/back chevron toggle; selection highlight not carried to the back view \\ \hline
\end{tabular}
\end{table}

% ============================================================
\section{Functional Test Cases}
% ============================================================

\begin{table}[H]
\centering
\caption{Functional test case results.}
\label{tab:functional_tests}
\begin{tabular}{|p{5.5cm}|p{4cm}|c|}
\hline
\textbf{Test Case} & \textbf{Expected Outcome} & \textbf{Result} \\ \hline
Patient registration and login & Account created; redirected to Home & Pass \\ \hline
Specialist registration and login & Specialist role unlocked & Pass \\ \hline
Capture image via camera & Image previewed & Pass \\ \hline
Upload image from gallery & Image previewed & Pass \\ \hline
Normal skin image submitted & Gate $P\le0.60$; ``No Disease Detected'' & Pass \\ \hline
Borderline image submitted & Gate gray zone; ``Inconclusive — retake'' & Pass \\ \hline
Diseased image (Acne/Eczema/Tinea) & Gate $\ge0.90$; CNN classifies; heatmap shown & Pass \\ \hline
Result saved to history & Entry appears with date, class, confidence & Pass \\ \hline
History entry deleted & Entry removed from SQLite & Pass \\ \hline
Airplane mode — full scan & Pipeline completes offline & Pass \\ \hline
Low-resolution input & Graceful upscaling; no crash & Pass \\ \hline
\end{tabular}
\end{table}

% ============================================================
\section{Boundary and Edge Case Testing}
% ============================================================

\begin{table}[H]
\centering
\caption{Boundary and edge case test results.}
\label{tab:edge_cases}
\begin{tabular}{|p{4cm}|p{4cm}|p{4cm}|c|}
\hline
\textbf{Input / Condition} & \textbf{Reason} & \textbf{Observed} & \textbf{Result} \\ \hline
50$\times$50\,px image & Below model input size & Upscaled; pipeline completed & Pass \\ \hline
12\,MP image (4032$\times$3024) & Typical phone output & Downscaled to $\le$1280\,px; no memory error & Pass \\ \hline
Non-skin photo (plain wall) & Unintended input & Gate scores low $P$; returns ``No Disease Detected'' rather than a confident disease label & Pass \\ \hline
Near-black underexposed image & Poor lighting & Completed; no crash & Pass \\ \hline
Overexposed near-white image & Bright outdoor capture & Completed; gate evaluated normally & Pass \\ \hline
Zero-byte corrupted file & Bad gallery selection & Decoder returned null; handled gracefully & Pass \\ \hline
5 consecutive scans & Repeated-use stress & No leak/crash/DB corruption & Pass \\ \hline
Rotation mid-analysis & Orientation change & Screen rebuilt; navigation completed & Pass \\ \hline
Camera permission denied & First-launch rejection & Explained and redirected to Settings & Pass \\ \hline
50+ history entries & Long-term usage & Loaded and scrolled smoothly & Pass \\ \hline
\end{tabular}
\end{table}

\noindent The replacement of the VAE with the supervised gate also improved the
non-skin case: a plain wall or fabric now scores a low $P(\text{disease})$ and is
correctly returned as ``No Disease Detected'' instead of being forwarded to the
CNN for a meaningless label.

% ============================================================
\section{On-Device Performance}
% ============================================================

\begin{table}[H]
\centering
\caption{On-device inference time per pipeline stage.}
\label{tab:inference_time}
\begin{tabular}{|l|c|c|}
\hline
\textbf{Stage} & \textbf{Mean (ms)} & \textbf{Std (ms)} \\ \hline
Image preprocessing (decode, EXIF bake, resize) & \textit{TODO} & \textit{TODO} \\ \hline
Stage 1 — Normal-vs-Disease gate (EfficientNetB0) & \textit{TODO} & \textit{TODO} \\ \hline
Stage 2 — CNN ensemble (B2+B3) & \textit{TODO} & \textit{TODO} \\ \hline
Stage 3 — Score-CAM heatmap & \textit{TODO} & \textit{TODO} \\ \hline
\textbf{Total end-to-end} & \textbf{\textit{TODO}} & \textit{TODO} \\ \hline
\end{tabular}
\end{table}

%% Fill from the device [TIMING] logcat lines (preprocess/gate/cnn/scorecam ms).

\begin{table}[H]
\centering
\caption{On-device model footprint (TFLite assets).}
\label{tab:storage}
\begin{tabular}{|l|c|}
\hline
\textbf{Component} & \textbf{Size (MB)} \\ \hline
Normal-vs-Disease gate (\texttt{normal\_gate.tflite}) & 4.6 \\ \hline
CNN B2 (\texttt{cnn\_b2\_model.tflite})                & 8.6 \\ \hline
CNN B3 (\texttt{cnn\_b3\_model.tflite})                & 11.9 \\ \hline
Score-CAM feature extractor (\texttt{b3\_feature\_extractor.tflite}) & 11.5 \\ \hline
\textbf{Total model assets} & \textbf{36.6} \\ \hline
\end{tabular}
\end{table}

% ============================================================
\section{Usability Testing}
% ============================================================

Three participants of differing backgrounds — a non-technical patient, a general
practitioner, and a university student — were each asked, without instruction,
to (i)~register and complete a scan, (ii)~find a past scan in History and view
its heatmap, and (iii)~add and save a clinical note.
Table~\ref{tab:usability} summarises the outcome.

\begin{table}[H]
\centering
\caption{Usability test results across three participants.}
\label{tab:usability}
\begin{tabular}{|l|c|c|p{5cm}|}
\hline
\textbf{Task} & \textbf{Avg.\ Time} & \textbf{Success} & \textbf{Key Observation} \\ \hline
Register and complete a scan & 1\,m\,42\,s & 3/3 & ``Tap To Scan'' found without guidance \\ \hline
Find past scan in History & 2\,m\,10\,s & 3/3 & Bottom navigation was intuitive \\ \hline
Add and save a clinical note & 3\,m\,05\,s & 2/3 & One user missed the notes field at first \\ \hline
\end{tabular}
\end{table}

\noindent The notes-field discoverability issue was addressed by adding
placeholder text (``Add a note for your doctor\ldots''); in a repeat evaluation
the field was noticed immediately. The GP confirmed the heatmap and confidence
score were presented as decision support rather than a replacement for clinical
judgement.

% ============================================================
\section{Use Case Scenarios}
% ============================================================

\subsection{Scenario 1: Patient Self-Check — Tinea}
The patient captures a circular rash. \textbf{Gate:} $P(\text{disease})$ is high
($\ge0.90$), so the image proceeds to the CNN. \textbf{CNN:} returns
\textbf{Tinea}. \textbf{Score-CAM:} highlights the annular scaly border. The
result, confidence, probability bars, and heatmap are saved to history.

\subsection{Scenario 2: Clinician Offline Decision Support — Eczema}
Offline, the GP photographs an inflamed lesion and is unsure between Eczema and
Tinea. \textbf{Gate:} $\ge0.90$ $\rightarrow$ CNN. \textbf{CNN:} returns
\textbf{Eczema}; \textbf{Score-CAM} highlights diffuse inflamed patches rather
than an active ring. Combined with examination, this supports the correct
treatment, and the GP saves a clinical note.

\subsection{Scenario 3: Patient Self-Check — Acne}
A student photographs facial breakouts. \textbf{Gate:} $\ge0.90$. \textbf{CNN:}
returns \textbf{Acne}; the heatmap activates over follicular clusters. The
timestamped result is saved for a future dermatology visit.

\subsection{Scenario 4: Normal Skin — Early Exit}
The patient photographs a faint mark. \textbf{Gate:} $P(\text{disease}) \le 0.60$,
classified \textbf{Normal}. The CNN and Score-CAM are not invoked, saving battery
and latency, and the screen shows ``No Disease Detected''.

\subsection{Scenario 5: Borderline — Abstention}
The patient submits a blurry, poorly framed photo. \textbf{Gate:}
$0.60 < P < 0.90$, so the app returns ``Inconclusive'' and asks for a clearer
retake instead of guessing — the safety behaviour the gray-zone band exists for.

\subsection{Scenario 6: Follow-Up via Scan History}
The patient compares two Eczema entries over two weeks; the more recent shows a
lower confidence and a less-activated heatmap, suggesting response to treatment,
and exports a PDF report for the dermatologist.

% ============================================================
\section{Chapter Summary}
% ============================================================

All 23 Dart unit tests pass, covering the model, security (PBKDF2 password
hashing), progression analytics, and the body-part selection UI. Functional,
boundary, and usability testing confirmed robust, crash-free behaviour across
extreme inputs and three user archetypes, and six use-case scenarios demonstrated
correct end-to-end behaviour including the normal early-exit and the borderline
abstention paths. On-device performance figures (per-stage latency) are captured
from the deployed application; the model footprint totals 36.6\,MB across the
four bundled TFLite assets, comfortably within a mobile storage budget for a
fully offline pipeline.
